\documentclass[10pt,a4paper]{amsart}
\usepackage[margin=19mm]{geometry}
\usepackage{amsmath,amssymb,mathtools}
\usepackage[scr=rsfs,cal=euler]{mathalfa}
\usepackage{xfrac}
\usepackage[unicode,hidelinks,bookmarks=false]{hyperref}
\newcommand{\ie}{i.e.\ }
\newcommand{\cf}{cf.\ }
\newcommand{\resp}{respectively\ }
\newcommand{\Subsection}{\subsection}
\providecommand{\nobreakdash}{\nobreak\hbox{-}}
\setlength{\emergencystretch}{2em}
\multlinegap0pt
\usepackage{bm}
\usepackage{bbm}
\usepackage{dsfont}
\usepackage{xspace}
\usepackage{cancel}
\usepackage{mathtools}
\usepackage{amsthm}
\makeatletter
\@ifundefined{thmcounter}{\newtheorem{thmcounter}{Thmcounter}}{}
\@ifundefined{lemma}{\newtheorem{lemma}[thmcounter]{Lemma}}{}
\@ifundefined{proposition}{\newtheorem{proposition}[thmcounter]{Proposition}}{}
\@ifundefined{remark}{\newtheorem{remark}[thmcounter]{Remark}}{}
\@ifundefined{theorem}{\newtheorem{theorem}[thmcounter]{Theorem}}{}
\makeatother
\makeatletter
\expandafter\ifx\csname EX\endcsname\relax
\newenvironment{EX}[1]{\vspace{.3cm}\begin{namelist}{000000}\item[#1]}{\end{namelist}}
\fi
\newcommand{\nP}{\mathbf{P}}
\expandafter\ifx\csname namelist\endcsname\relax
\newenvironment{namelist}[1]{%
\begin{list}{}
{\renewcommand{\makelabel}{\namelistlabel}
\settowidth{\labelwidth}{#1}%
\setlength{\labelsep}{0.3em}%
\setlength{\leftmargin}{\labelwidth}%
\addtolength{\leftmargin}{\labelsep}}}{%
\end{list}}
\fi
\newcommand{\sA}{\mathscr{A}}
\newcommand{\sF}{\mathscr{F}}
\newcommand{\sO}{\mathscr{O}}
\makeatother
\title[Ulrich sheaves and determinantal representations for higher ]{Ulrich sheaves and determinantal representations for higher secant varieties of curves}
\author{Daniele Agostini, Mario Kummer, Jinhyung Park}
\date{Source: arXiv:2607.22247v1 (CC BY 4.0)}
\begin{document}
\maketitle

\noindent\emph{Source and licence.} Ulrich sheaves and determinantal representations for higher secant varieties of curves. Version arXiv:2607.22247v1 (CC BY 4.0), distributed under the Creative Commons Attribution 4.0 International licence, as recorded in the arXiv licence metadata for this exact version and shown on the abstract page https://arxiv.org/abs/2607.22247v1. Full rights notice: licenses/arxiv-2607.22247-CC-BY-4.0.txt.\par\smallskip

\noindent\textit{Editorial scope.} Counted unit: the Theorem whose source label is \texttt{thm:detrep} in the article (source0.tex lines 1231--1235, source label \texttt{thm:detrep}), delivered together with the article's own complete original proof of it (source0.tex lines 1236--1279, one \texttt{proof} environment of 44 lines). No mathematics is written, repaired or re-derived: statement and proof are the article's own text line for line. Required context retained uncounted: thm:admrepresentation (lines 634--647); prop:szego (lines 1051--1056); rem:szegosymm (lines 1086--1088); lem:phivanishes (lines 1216--1224). Every definition, display and label of those lines is retained unchanged. Omitted from this package and not redistributed: 1--633, 648--1050, 1057--1085, 1089--1215, 1225--1230, 1280--2189. Nothing inside a retained excerpt is omitted or altered: every retained body is delivered exactly as the article's lines are written, so no omission receipt of any kind accompanies this package. Citations: the retained excerpts carry no citation command at all, so no bibliography entry of the article is transcribed and no reference is redistributed. Reference adaptation: none was needed and none was made; every reference inside the delivered unit resolves inside it, so no unresolved reference or citation is produced. Printed slips kept literal: no printed word, symbol, hypothesis or number of the counted statement or of its proof was altered. Layout and class: the article's own class is \texttt{amsart} with packages including geometry, amscd, amssymb, amsmath, wasysym, graphicx, amsfonts, mathrsfs, eucal, latexsym, verbatim, xy, xcolor, bookmark; the wrapper is the lane's pinned \texttt{amsart} editorial wrapper at 10pt on A4, which restates the theorem-like environments the retained text needs and adds the article's own macro definitions that the retained text needs (\texttt{\textbackslash nP}, \texttt{\textbackslash sA}, \texttt{\textbackslash sF}, \texttt{\textbackslash sO}), with extra wrapper packages (bm, bbm, dsfont, xspace, cancel, mathtools). The delivered numbering is the unit's own. Assets: no retained span contains a figure environment, a graphics-inclusion command, a PDF-page inclusion, a table or a TikZ picture; no image file is copied into this package, the package declares no asset dependency and no image file is redistributed with it. Licence: version arXiv:2607.22247v1 of this article is distributed under the Creative Commons Attribution 4.0 International licence, as recorded in the arXiv licence metadata for this exact version and shown on the abstract page https://arxiv.org/abs/2607.22247v1. A full rights notice is delivered at licenses/arxiv-2607.22247-CC-BY-4.0.txt. Characters: every retained excerpt is pure ASCII, so the delivered engine, XeLaTeX with its default Latin Modern text font, reports no missing character.
\begin{theorem}\label{thm:admrepresentation} 
Let $X\subseteq \nP^r$ be an irreducible nondegenerate projective variety of dimension $n$ as before and let $\sF,\sF'$ two rank one and torsion-free sheaves on $X$ with $h^0(X,\sF)=h^0(X,\sF')=d = \deg(X)$. Consider a linear map
\[ \gamma\colon \wedge^{n+1}V \longrightarrow H^0(X\times X,\sF\boxtimes \sF') = H^0(X,\sF)\otimes H^0(X,\sF')\]
with the following properties:
\begin{itemize}
	\item[(i)] If $p_0,\dots,p_n \in V$ have no common zero on $X$, then a general fiber of the projection $\pi = (p_0:\dots :p_n)\colon X\to \nP^n$ consists of $d$ distinct points $\{x_1,\dots,x_d\}$ such that
	\[ \gamma(p_0\wedge \dots \wedge p_n)(x_i,x_j) = \begin{cases} \text{ nonzero} & \text{if } i=j \\ \text{ zero} &\text{if } i\ne j \end{cases}.\]
	\item[(ii)] If $x\in X$ is a general point and if  $p_0,\dots,p_n\in V$ are general elements vanishing at $x$, then
	\[ \gamma(p_0\wedge \dots \wedge p_n)(x,-)\footnote{Here we see the restriction $\gamma(p_0\wedge \dots \wedge p_n)(x,-)$ to $\{x\}\times X$ as a section of $H^0(X,\sF')$} = 0 \quad \text{ in } H^0(X,\sF'). \]
	Furthermore,  a general fiber of the projection $\pi = (p_0:\dots:p_n)\colon X\dashrightarrow \nP^{n}$ consists of $d-1$ points $\{x_1,\dots,x_{d-1}\}$ such that
	\[ \gamma(p_0\wedge \dots \wedge p_n)(x_i,x_j) = \begin{cases} \text{ nonzero} & \text{if } i=j \\ \text{ zero} &\text{if } i\ne j \end{cases}.\]
\end{itemize}
Then $\gamma$ is an admissible determinantal representation of degree $d$ of $X$, and the corresponding Ulrich sheaf is the image of the evaluation map $H^0(X,\sF)\otimes \sO_X \to \sF$. In particular it coincides with $\sF$ if and only if $\sF$ is globally generated.
\end{theorem}

\begin{proposition}\label{prop:szego}
	With the above notation, there is a unique $(k+1)$-Szeg\H{o} kernel $S_{(\alpha,\alpha')}$ associated to $(\alpha,\alpha')$. Furthermore
	\[ S_{(\alpha,\alpha')}(\xi,\xi')= 0 \qquad  \text{ if and only if } \qquad  h^0(C,\alpha(\xi'-\xi)) = h^0(C,\alpha'(\xi-\xi')) >0.   \]
	and under the involution of $C_{k+1}\times C_{k+1}$ that exchanges the factors, it holds that
	\[ S_{(\alpha,\alpha')}(\xi,\xi') = S_{(\alpha',\alpha)}(\xi',\xi) \qquad  \text{ for all } \xi,\xi'\in C_{k+1}. \]
\end{proposition}

\begin{remark}\label{rem:szegosymm}
 If $\alpha$ is a theta characteristic with no cohomology, Proposition \ref{prop:szego} shows that $S_{(\alpha,\alpha)}$ is symmetric: $S_{(\alpha,\alpha)}(\xi,\xi') = S_{(\alpha,\alpha)}(\xi',\xi)$ for all $\xi,\xi'\in C_{k+1}$.
\end{remark}

\begin{lemma}\label{lem:phivanishes}
	Let $p_0,\dots,p_{2k+1} \in V$ and let $\langle p_0,\dots,p_{2k+1} \rangle \subseteq V$ be the subspace that they generate. For $\xi,\xi'\in C_{k+1}$ it holds that
	\[ \phi(p_0\wedge \dots \wedge p_{2k+1})(\xi,\xi') = 0 \]
	if and only if one or both of the following two condition hold:
	\begin{itemize}
		\item[$(a)$]  $\langle p_0,\dots,p_{2k+1} \rangle$ does not separate $\xi+\xi'$.
		\item[$(b)$] $h^0(C,\alpha(\xi'-\xi))>0$.
	\end{itemize}
\end{lemma}

\begin{theorem}\label{thm:detrep}
	Assume that $r\geq 2k+2$ and that \[d:=h^0(C_{k+1},A_{k+1,L\otimes \alpha}) = h^0(C_{k+1},A_{k+1,L\otimes \alpha'}) = \deg(\Sigma_k).\] The linear map
	\[  \phi = \phi_{k,V,\alpha} \colon \wedge^{2k+2} V \longrightarrow H^0(C_{k+1},A_{k+1,L\otimes \alpha}) \otimes H^0(C_{k+1},A_{k+1,L\otimes \alpha'}) \]
	is an admissible determinantal representation of $\Sigma_k$ of degree $d$, whose corresponding Ulrich sheaf , that we denote by $\sA_{k,L\otimes \alpha}$, has rank one. If $\alpha$ is a theta-characteristic, this is a symmetric representation and $\sA_{k,L\otimes \alpha}$ is a symmetric Ulrich sheaf.
\end{theorem}

\begin{proof}
	We first treat the case when $V$ is $k$-very ample so that there is  the well-defined morphism $\beta_k\colon B_k \longrightarrow \Sigma_k$ and the  coherent sheaf $\widetilde{\sA}_{k,L\otimes \alpha} = \beta_{k,*}(\pi_k^* A_{k+1,L\otimes \alpha})$ on $\Sigma_k$. By construction, $H^0(\Sigma_k,\widetilde{\sA}_{k,L\otimes \alpha}) \cong H^0(B_k,\pi_k^*A_{k+1,L\otimes \alpha}) \cong H^0(C_{k+1},A_{k+1,L\otimes \alpha})$. Furthermore, if $x \in \Sigma_k$ is a general point, then the fiber $\beta^{-1}(x)$ is a single point, meaning that there is a unique $\xi\in C_{k+1}$ such that $x\in \langle \xi \rangle$. The evaluation of a section of $H^0(\Sigma_k,\widetilde{\sA}_{k,L\otimes \alpha})$ at $x$ corresponds to the evaluation of a section of $H^0(C_{k+1},A_{k+1,L\otimes \alpha})$ at $\xi$. To apply Theorem \ref{thm:admrepresentation}, we first observe that $\widetilde{\sA}_{k,L\otimes \alpha}$ is torsion-free and of rank one because $\beta_k$ is birational.  Now we check the two properties:
	\medskip

	\noindent
	(i)  Suppose that $p_0, \ldots, p_{2k+1}$ have no common zero on $\Sigma_k$ and let
	$
	\pi  \colon \Sigma_k \longrightarrow \nP^{2k+1},
	$
	be the corresponding outer projection,
	which is a finite morphism of degree $d$. Consider a general fiber $\{y_1, \ldots, y_d\}$. We can assume that there are unique $\xi_1, \ldots, \xi_d \in C_{k+1}$ such that $y_i \in \langle \xi_i \rangle$ for all $1 \leq i \leq d$. We may also assume that each $\xi_i$ consists of $k+1$ distinct points and that $\pi$ is an immersion around $y_1, \ldots, y_d$. We want to show that
	$$
	\phi(p_0 \wedge \cdots \wedge p_{2k+1})(\xi_i, \xi_j) =
	\begin{cases} \text{nonzero} & \text{if $i=j$} \\ \text{zero} & \text{if $i \neq j$} \end{cases}.
	$$
	First, assume that $i \neq j$. Note that $\pi$ sends $\langle \xi_i + \xi_j \rangle$ to a smaller dimensional linear subspace of $\nP^{2k+1}$. This means that $\langle p_0,\dots,p_{2k+1}\rangle$ does not separate $\xi_i+\xi_j$. By Lemma \ref{lem:phivanishes} we get $\phi(p_0 \wedge \cdots \wedge p_{2k+1})(\xi_i, \xi_j) = 0$. Next, assume that $i=j$. Terracini's lemma tells us that the projective tangent space to $\Sigma_{k}$ at $y_i$ is $\langle 2\xi_i \rangle$. As $\pi$ is an immersion at $x_i$, this means that $\langle p_0,\dots,p_{2k+1} \rangle$ separates $2\xi_i$. By Lemma \ref{lem:phivanishes} it is then enough to check that $S_{(\alpha,\alpha')}(\xi_i,\xi_i)\ne 0$ and this follows from Proposition \ref{prop:szego}.
	\medskip

	\noindent
	(ii) Suppose that $p_0, \ldots, p_{2k+1}$ have a common zero at a point $y$ on $\Sigma_k$. There is $\xi \in C_{k+1}$ such that $y\in \langle \xi \rangle$. This means that $\langle p_0,\dots,p_{2k+1} \rangle$ does not separate $\xi$, so that it does not separate $\xi+\xi'$ either, for any  $\xi'\in C_{k+1}$. Then Lemma \ref{lem:phivanishes} shows that
	$$
	\phi(p_0 \wedge \cdots \wedge p_{2k+1})(\xi, \xi')=0~~\text{ for any $\xi' \in C_{k+1}$}.
	$$
	Furthermore, if $x\in X$ is a general point and if $p_0,\dots,p_{2k+1}$ have a unique common zero $x$ on $\Sigma_k$, then they induce an inner projection $\pi\colon \Sigma_k \dashrightarrow \nP^{2k+1}$ of degree $d-1$. If $\{y_1,\dots,y_{d-1}\}$ is a general fiber, and if $\xi_1,\dots,\xi_{d-1} \in C_{k+1}$ are the unique schemes such that $y_i \in \langle \xi_i \rangle$, then the same arguments used for point (i) show that
	$$
	\phi(p_0 \wedge \cdots \wedge p_{2k+1})(\xi_i, \xi_j) =
	\begin{cases} \text{nonzero} & \text{if $i=j$} \\ \text{zero} & \text{if $i \neq j$} \end{cases}.
	$$
	which is what we needed to prove.
	\medskip

	\noindent
	The previous discussion shows that we can apply Theorem \ref{thm:admrepresentation}, so that $\phi$ induces an admissible determinantal representation of degree $d$ of $\Sigma_d$, whose corresponding Ulrich sheaf $\sA_{k,L\otimes \alpha}$ is the image of the evaluation map
	\[  H^0(C_{k+1},A_{k+1,L\otimes \alpha}) \otimes \sO_{\Sigma_k} \longrightarrow \widetilde{\sA}_{k,L\otimes \alpha}. \]
	\medskip

	\noindent
	Assume now that $V$ is not necessarily $k$-very ample. Then the map $\beta_k\colon B_k \dashrightarrow \Sigma_k$ is only rational, but it can be resolved via two birational projective maps $\varepsilon_k \colon \widetilde{B}_k \to B_k$ and $\widetilde{\beta}_k\colon \widetilde{B}_k \to \Sigma_k$, such that $\widetilde{\beta}_k = \beta_k \circ \varepsilon_k$. Notice that  $\varepsilon_{k,*}\sO_{\widetilde{B_k}} = \sO_{B_k}$, so that  $H^0(\widetilde{B}_k,\varepsilon_k^*\pi_k^*A_{k+1,L\otimes \alpha}) \cong H^0(B_k,\pi_k^*A_{k+1,L\otimes \alpha}) \cong H^0(C_{k+1},A_{k+1,L\otimes \alpha})$. Consider now the coherent sheaf $\widetilde{\sA}_{k,L\otimes \alpha} = \widetilde{\beta}_{k,*}(\varepsilon_k^* \pi_k^* A_{k+1,L\otimes \alpha})$ on $\Sigma_k$. We see that $H^0(\Sigma_k,\widetilde{\sA}_{k,L\otimes \alpha}) \cong H^0(C_{k+1},A_{k+1,L\otimes \alpha})$, and then we can use the same proof of the case when $V$ is $k$-very ample.
	\medskip


	\noindent
	Finally, if $\alpha$ is a theta characteristic, Remark \ref{rem:szegosymm}, shows that the Szeg\H{o} kernel $S_{(\alpha,\alpha)}$ is symmetric. Furthermore, the addition map $\sigma\colon C_{k+1}\times C_{k+1} \to C_{2k+2}$ is also symmetric, so that for any $p_0,\dots,p_{2k+1}\in V$ we see that \[ \phi(p_0\wedge \dots \wedge p_{2k+1}) = \sigma^*\det_{2k+2,L}\left( p_0 \wedge \dots \wedge p_{2k+1}  \right)\cdot S_{(\alpha,\alpha)} \]  is a symmetric element of $H^0(C_{k+1},A_{k+1,L\otimes \alpha}) \otimes H^0(C_{k+1},A_{k+1,L\otimes \alpha})$. This proves that $\phi$ yields a symmetric determinantal representation.
\end{proof}

\end{document}
